% --- Содержимое файла materials/5_1.tex ---

\subsection{Линейные операторы: ограниченность и непрерывность}

Пусть $E_1, E_2$ --- линейные нормированные пространства над полем $\mathbb{K}$ ($\mathbb{R}$ или $\mathbb{C}$).
Отображение $A: E_1 \to E_2$ будем называть \textbf{оператором}, а отображение $f: E \to \mathbb{K}$ --- \textbf{функционалом}.

\begin{definition}[Ограниченный оператор]
	Оператор $A$ называется \textbf{ограниченным}, если он переводит ограниченные множества в ограниченные.
	
	Формально: для любого ограниченного множества $M \subset E_1$ образ $A(M)$ ограничен в $E_2$.
\end{definition}


\begin{idea}
	\textbf{Интуиция:} Ограниченный линейный оператор — это такой оператор, который «не растягивает» векторы бесконечно сильно. У него есть конечный «коэффициент усиления». Если вы подаете на вход вектор длины 1, на выходе не может получиться вектор бесконечной длины.
\end{idea}

\begin{definition}[Образ и Ядро]
	Для линейных операторов вводятся следующие множества:
	\begin{itemize}
		\item \textbf{Образ оператора} (Image):
		\[ \operatorname{Im} A = \{ y \in E_2 \mid \exists x \in E_1 : Ax = y\} \]
		\item \textbf{Ядро оператора} (Kernel):
		\[ \operatorname{Ker} A = \{ x \in E_1 \mid Ax = 0\} \]
	\end{itemize}
\end{definition}


\begin{definition}[Непрерывный оператор]
	Линейный оператор $A$ называется \textbf{непрерывным}, если для любой сходящейся последовательности $x_n \to x$ выполнено $Ax_n \to Ax$.
\end{definition}


\begin{theorem}[Критерий непрерывности линейного оператора]\label{thm:bounded_cont}
	Пусть $E_1, E_2$ --- линейные нормированные пространства, $A: E_1 \to E_2$ --- линейный оператор.
	
	Тогда $A$ ограничен $\iff$ $A$ непрерывен.
\end{theorem}

\begin{proof}
	\textbf{1. Ограниченность $\Rightarrow$ Непрерывность.}
	
	Пусть $A$ ограничен. Воспользуемся эквивалентным определением ограниченности (через норму оператора, см. утверждение \ref{equiv-bounded-oper}): $\exists C > 0: \|Ax\| \leqslant C\|x\|$.
	
	Пусть $x_n \to x$. Это равносильно тому, что $\|x_n - x\|_{E_1} \to 0$. Оценим норму разности образов, используя линейность $A$:
	\[
	\|Ax_n - Ax\|_{E_2} = \|A(x_n - x)\|_{E_2} \leqslant \|A\| \cdot \|x_n - x\|_{E_1}
	\]
	Так как $\|x_n - x\|_{E_1} \to 0$, то и $\|Ax_n - Ax\|_{E_2} \to 0$, что означает $Ax_n \to Ax$.
	
	\textbf{2. Непрерывность $\Rightarrow$ Ограниченность.}
	
	Докажем от противного. Предположим, что $A$ непрерывен, но \textbf{не} является ограниченным.
	
	Это значит, что его «коэффициент растяжения» может быть сколь угодно большим:
	\[ \forall n \in \mathbb{N} \quad \exists x_n \in E_1 : \|Ax_n\|_{E_2} > n \|x_n\|_{E_1} \]
	Очевидно, $x_n \neq 0$. Рассмотрим нормированную и масштабированную последовательность векторов:
	\[ y_n = \frac{x_n}{n \cdot \|x_n\|_{E_1}} \]
	
	Заметим два факта:
	\begin{enumerate}
		\item Последовательность $y_n$ стремится к нулю:
		\[ \|y_n\|_{E_1} = \frac{\|x_n\|}{n \|x_n\|} = \frac{1}{n} \to 0 \implies y_n \to 0 \]
		\item Из непрерывности оператора $A$ следует, что образ должен стремиться к образу нуля:
		\[ A y_n \to A(0) = 0 \]
	\end{enumerate}
	
	Однако, оценим норму образа $y_n$ снизу:
	\[ \|A y_n\|_{E_2} = \left\| A \left( \frac{x_n}{n \|x_n\|} \right) \right\| = \frac{1}{n \|x_n\|} \|A x_n\| > \frac{1}{n \|x_n\|} \cdot (n \|x_n\|) = 1 \]
	
	Получили противоречие: последовательность $A y_n$ должна стремиться к $0$, но норма всех ее элементов больше $1$. Следовательно, предположение неверно, и оператор $A$ ограничен.
\end{proof}

\begin{idea}[Смысл теоремы]
	
	Для линейных операторов понятия «не порвать график» (непрерывность) и «не растягивать бесконечно сильно» (ограниченность) совпадают.
	
	Это очень сильное свойство, которого нет у обычных функций (например, $f(x) = x^2$ непрерывна, но не ограничена на $\mathbb{R}$). 
	
	Причина в \textbf{линейности}: оператор одинаково устроен во всем пространстве. Если он «ведет себя хорошо» (непрерывен) в одной точке (в нуле), он обязан быть ограниченным везде. И наоборот: если у него есть конечный «коэффициент растяжения», он не может совершать скачков.
	
	
	
	2. Идея доказательства 
	
	[Логика доказательства]
	
	\textbf{1. Ограниченность $\Rightarrow$ Непрерывность:}
	
	Здесь работает простая оценка. Если оператор растягивает любой вектор максимум в $\|A\|$ раз, то маленькая разность на входе ($x_n - x \to 0$) превратится в маленькую разность на выходе ($A(x_n - x) \to 0$).
	
	\textbf{2. Непрерывность $\Rightarrow$ Ограниченность (от противного):}
	
	Если оператор \textit{неограничен}, он работает как «усилитель» с бесконечной мощностью: найдутся векторы, которые он удлиняет в $100, 1000, n$ раз.
	
	Мы используем линейность, чтобы «обмануть» его: берем эти векторы и уменьшаем их длину до микроскопических размеров ($y_n \to 0$). Но так как коэффициент растяжения растет быстрее ($n$), на выходе мы все равно получаем «большие» векторы ($\|Ay_n\| > 1$). Это противоречит непрерывности в нуле.
	
\end{idea}